% The sample deck offered on the opening page: builds, speaker notes, an outline, links between
% slides, and TikZ pictures that animate. Build with `make` in this directory; the PDF is committed.
\documentclass[aspectratio=169,14pt]{beamer}
\usetheme{default}
\usepackage{pdfpc}
\usepackage{pgfplots}
\pgfplotsset{compat=1.17}
\usetikzlibrary{arrows.meta,positioning}
\definecolor{lectern}{HTML}{0736A8}
\definecolor{warm}{HTML}{C2410C}
\setbeamercolor{structure}{fg=lectern}
\setbeamercolor{frametitle}{fg=lectern}
\setbeamerfont{frametitle}{series=\bfseries}
\setbeamertemplate{navigation symbols}{}
\setbeamerfont{subtitle}{size=\normalsize,series=\mdseries}
\setbeamerfont{title}{size=\huge,series=\bfseries}
\setbeamertemplate{itemize items}[circle]
\setbeamertemplate{footline}{\hfill\usebeamercolor[fg]{structure}\footnotesize\insertframenumber\hspace{1em}\vspace{0.6em}}
\tikzset{
  box/.style={draw=lectern, very thick, rounded corners=3pt, minimum width=2.6cm, minimum height=1.1cm, align=center},
  flow/.style={-{Stealth[length=7pt]}, very thick, lectern},
}
\title{A tour of Lectern}
\subtitle{A sample deck with builds, notes, links and animation}
\author{Made with Beamer and TikZ}
\date{}
\begin{document}

\begin{frame}
  \titlepage
  \pdfpcnote{Welcome. These are speaker notes: only the presenter view shows them, never the audience.\\\\Press the right arrow, or click the slide, to go on.}
\end{frame}

\section{Moving around}

\begin{frame}{Frames and builds}
  A slide that appears step by step is one \emph{frame} with several \emph{builds}.
  \begin{itemize}
    \item<2-> \textbf{Arrow keys} step one build at a time.
    \item<3-> \textbf{Shift with an arrow} moves a whole frame.
    \item<4-> \textbf{A number, then Enter,} goes to that frame.
    \item<5-> \textbf{Tab} shows every frame at once.
  \end{itemize}
  \only<1>{\pdfpcnote{This frame has five builds. The frame counter stays at 2 while they appear.}}
  \only<2>{\pdfpcnote{Each build can carry a note of its own. This one belongs to the build about the arrow keys.}}
  \only<3>{\pdfpcnote{Try it: Shift and the right arrow skips the remaining builds.}}
  \only<4>{\pdfpcnote{Type 5 and press Enter to go to frame 5. Delete or Backspace returns here.}}
  \only<5>{\pdfpcnote{The overview keeps working while you present: the keys still move the talk.}}
\end{frame}

\begin{frame}[label=links]{Links between slides}
  Links in the PDF work with a click, in either view.

  \bigskip
  \hyperlink{backup-proof}{\scalebox{1.5}{\beamergotobutton{The proof, in the backup slides}}}

  \medskip
  \hyperlink{backup-data}{\scalebox{1.5}{\beamergotobutton{The data behind the picture}}}

  \bigskip
  After a jump, \textbf{Back} (or Delete) returns to where you were.
  \pdfpcnote{Both buttons jump to the appendix. Each backup slide has a return button, and the Back button under the current slide does the same.\\\\- The outline tab lists the sections.}
\end{frame}

\section{Pictures that move}

\begin{frame}{A curve drawn point by point}
  \centering
  \begin{tikzpicture}
    \begin{axis}[width=12.5cm, height=6.4cm, axis lines=left, xmin=0, xmax=1, ymin=0, ymax=0.3,
        xlabel={Tax rate}, ylabel={Revenue}, xtick={0,0.25,0.5,0.75,1}, ytick=\empty,
        every axis plot/.append style={line width=1.6pt}, clip=false]
      \addplot[lectern!25, domain=0:1, samples=60] {x*(1-x)};
      \foreach \n in {1,...,20} {
        \only<\n>{
          \pgfmathsetmacro{\t}{\n/20}
          \addplot[lectern, domain=0:\t, samples=60] {x*(1-x)};
          \addplot[only marks, mark=*, mark size=3.5pt, warm] coordinates {(\t, {\t*(1-\t)})};
          \addplot[warm, dashed, line width=0.8pt] coordinates {(\t,0) (\t, {\t*(1-\t)})};
        }
      }
      \only<10>{\node[above=6pt, warm] at (axis cs:0.5,0.25) {the peak};}
    \end{axis}
  \end{tikzpicture}
  \pdfpcnote{Twenty builds, one per position of the point. Hold the right arrow to play it as an animation; the left arrow runs it backwards.\\\\Revenue is the rate times what is left of the base, so it peaks in the middle.}
\end{frame}

\begin{frame}{A diagram that assembles itself}
  \centering
  \begin{tikzpicture}[node distance=1.2cm and 0.9cm, font=\small]
    \node[box] (pdf) {Your PDF};
    \uncover<2->{\node[box, right=of pdf] (app) {Lectern,\\in a browser};}
    \uncover<3->{\node[box, above right=of app] (aud) {Slides for\\the audience};}
    \uncover<4->{\node[box, below right=of app] (pres) {Presenter view\\with notes};}
    \uncover<5->{\node[box, right=of pres, draw=warm] (phone) {A phone\\as a remote};}
    \uncover<2->{\draw[flow] (pdf) -- (app);}
    \uncover<3->{\draw[flow] (app) |- (aud);}
    \uncover<4->{\draw[flow] (app) |- (pres);}
    \uncover<5->{\draw[flow, warm, {Stealth[length=7pt]}-{Stealth[length=7pt]}] (pres) -- (phone);}
  \end{tikzpicture}
  \only<1-4>{\pdfpcnote{The picture gains a piece with each build. The next-slide preview shows what the next key press adds.}}
  \only<5>{\pdfpcnote{Any device that enters the talk code joins in, and each one chooses what it shows.}}
\end{frame}

\section{Tools while you talk}

\begin{frame}{Things to try}
  \begin{columns}[T]
    \begin{column}{0.48\textwidth}
      \textbf{Point and draw}
      \begin{itemize}
        \item \textbf{L} for a laser pointer
        \item \textbf{D} to draw on the slide
        \item \textbf{C} to clear the ink
      \end{itemize}
    \end{column}
    \begin{column}{0.48\textwidth}
      \textbf{Pause the room}
      \begin{itemize}
        \item \textbf{B} for a black screen
        \item \textbf{W} for a white one
        \item \textbf{?} for every shortcut
      \end{itemize}
    \end{column}
  \end{columns}
  \pdfpcnote{Notes understand a little Markdown: *emphasis*, `code`, and lists.\\\\- A- and A+ change the size of this text.\\- The clock above starts with the talk; T pauses it.}
\end{frame}

\begin{frame}{Now open your own}
  \begin{itemize}
    \item Any PDF works: Beamer, PowerPoint or Keynote exports.
    \item It opens in this browser and is not uploaded anywhere.
    \item Press \textbf{Esc} twice to go back and choose a file.
  \end{itemize}
  \pdfpcnote{The last frame of the main talk. What follows is the appendix, reached by the links on frame 3.}
\end{frame}

\appendix
\section{Backup slides}

\begin{frame}[label=backup-proof]{Backup: the proof}
  Revenue is $R(t) = t\,(1-t)$, so
  \[ R'(t) = 1 - 2t = 0 \quad\Longrightarrow\quad t^* = \tfrac12 . \]
  \uncover<2->{Since $R''(t) = -2 < 0$, this is the maximum, with $R(t^*) = \tfrac14$.}

  \bigskip
  \hyperlink{links}{\scalebox{1.5}{\beamerreturnbutton{Back to the links}}}
  \pdfpcnote{A backup slide with two builds. The return button, the Back button and the Delete key all lead back.}
\end{frame}

\begin{frame}[label=backup-data]{Backup: the data}
  \centering
  \begin{tikzpicture}
    \begin{axis}[width=11cm, height=5.6cm, ybar, bar width=14pt, axis lines=left, ymin=0, ymax=0.3,
        xlabel={Tax rate}, ylabel={Revenue}, xtick=data, ytick=\empty, enlarge x limits=0.12]
      \addplot[fill=lectern!70, draw=lectern] coordinates {(0.1,0.09) (0.3,0.21) (0.5,0.25) (0.7,0.21) (0.9,0.09)};
    \end{axis}
  \end{tikzpicture}

  \hyperlink{links}{\scalebox{1.5}{\beamerreturnbutton{Back to the links}}}
  \pdfpcnote{The same curve at five rates.}
\end{frame}

\end{document}
